\section{Results}

For running the application the following commands were used for compiling,
uploading the project and monitoring the port to observe the periodic reports.

\begin{verbatim}
# list connected boards and their ports
pio device list
# compile and upload to ESP32
pio run --target upload
# open serial monitor at 115200 baud
pio device monitor --baud 115200
\end{verbatim}

Below is the startup banner and sample serial output of the system:

\begin{verbatim}
==========================================
  Lab 7 - Actuator Control System (Var C)
==========================================
Analog  : L298 PWM on GPIO18 (LED test)
          IN1=GPIO16  IN2=GPIO17
Binary  : Relay on GPIO19 (active-HIGH)
Pipeline: Sat -> Median(5) -> EMA(a=0.30) -> Ramp(2/cyc)
Overload: >= 95%
Commands: S<0-100> DF DR X RON ROFF ?
==========================================

[CMD] Speed target -> 50%
>target:50,cond:8.0,current:4.0,relay:0
>target:50,cond:18.4,current:10.0,relay:0
>target:50,cond:26.9,current:16.0,relay:0

===== ACTUATOR REPORT (t=5000 ms) =====
Target speed  : 50%
Conditioned   : 44.2%  (median=5, EMA alpha=0.30)
Current speed : 40.0%  (ramp step=2/cyc)
Direction     : FORWARD
Motor state   : RAMP_UP
Relay         : OFF
Alert         : NONE
PWM duty      : 102/255
======================================

[CMD] Speed target -> 100%
>target:100,cond:95.0,current:80.0,relay:0
[CMD] Relay -> ON
>target:100,cond:98.0,current:82.0,relay:1

===== ACTUATOR REPORT (t=10000 ms) =====
Target speed  : 100%
Conditioned   : 99.1%  (median=5, EMA alpha=0.30)
Current speed : 88.0%  (ramp step=2/cyc)
Direction     : FORWARD
Motor state   : RAMP_UP
Relay         : ON
Alert         : OVERLOAD
PWM duty      : 224/255
======================================

[CMD] EMERGENCY STOP
>target:0,cond:0.0,current:0.0,relay:1
\end{verbatim}

\textbf{Hardware note:} The photographs in this section show an LED connected to
GPIO~18 (ENA) in place of a DC motor, as the motor was not available at the time the
photos were taken. The LED brightness is proportional to the PWM duty cycle and
visually confirms the ramping, smoothing, and speed control behavior. During the
actual laboratory presentation, however, a real \textbf{DC motor connected through
the L298N H-bridge driver} was used. The motor responded correctly to all speed,
direction, and emergency stop commands, demonstrating the full analog actuator
control pipeline on real hardware. No code changes were required to switch from the
LED to the motor, confirming the driver abstraction design.

The system behavior is as follows:
\begin{itemize}
  \item At startup, the OLED shows \texttt{-- Actuator Ctrl --} / \texttt{Lab7 Variant C}
    / \texttt{Starting...} for 1\,second while all modules initialize.
  \item After a speed command (e.g., \texttt{S50}), \texttt{currentSpeed} increments
    by $2\,\%$ every 50\,ms, with the OLED showing the motor state transition from
    IDLE $\to$ RAMP\_UP $\to$ RUNNING. The LED on GPIO~18 brightens gradually,
    visually confirming the ramp behavior in place of a real motor.
  \item The \textbf{conditioned speed} lags the target due to the median and EMA filters.
    Rapid command changes (e.g., \texttt{S10} then \texttt{S90} then \texttt{S30})
    are smoothed, and the ramp ensures the LED brightness changes at a controlled rate.
  \item Sending \texttt{S100} causes the conditioned speed to reach $\geq 95\,\%$,
    triggering \textbf{OVERLOAD} alert: the red LED turns on, the green LED turns off,
    and the OLED shows \texttt{Alrt:OVERLD}.
  \item Sending \texttt{S0} causes the motor state to transition to RAMP\_DOWN and
    eventually IDLE. Sending \texttt{X} instead causes an \textbf{immediate} stop
    (no ramping), resetting all filter state.
  \item \texttt{DR} followed by \texttt{S50} swaps the IN1/IN2 direction pins --
    the motor would reverse; the LED brightness is unchanged.
  \item \texttt{RON} energizes the relay (GPIO~19 HIGH) and the OLED shows
    \texttt{Relay: ON}. \texttt{ROFF} de-energizes it.
  \item \texttt{?} prints a full status dump to the serial terminal at any time.
  \item The serial plotter line \texttt{>target:X,cond:Y,current:Z,relay:R} is emitted
    every 500\,ms. When visualized in the VS Code Serial Plotter extension, the three
    speed traces show the filtering and ramp effects clearly: \texttt{target} steps
    instantly, \texttt{cond} follows with smoothing lag, and \texttt{current} ramps
    at a constant slope.
  \item Entering \texttt{S150} is rejected with an error message and clamped to
    $100\,\%$ before entering the pipeline.
\end{itemize}

\clearpage
Below is the visual representation of the circuit:

\begin{figure}[h]
    \centering
    \includegraphics[width=0.6\textwidth]{img/ledoff.jpg}
    \caption{Circuit during development -- LED used as analog actuator substitute}
    \label{fig:ledoff}
\end{figure}

Figure \ref{fig:ledoff} shows the circuit during the development phase, with an LED
on GPIO~18 acting as a visual substitute for the DC motor. Note that during the actual
laboratory presentation the LED was replaced by a real DC motor connected through the
L298N H-bridge driver.

\begin{figure}[h]
    \centering
    \includegraphics[width=0.6\textwidth]{img/ledon.jpg}
    \caption{Circuit at maximum output -- relay ON and full PWM duty (LED substitute)}
    \label{fig:ledon}
\end{figure}

Figure \ref{fig:ledon} shows the circuit at maximum speed command (\texttt{S100})
with the relay energized. The PWM LED is at full brightness, representing 100\% duty
cycle. During the presentation this same state drove the DC motor at full speed through
the L298N driver.
